\section{Commercial policy evaluation: what an offline score permits}
\label{exec16:commercial-policy}

A servicing manager comparing a contact policy with the existing workflow
needs a value for a specified action rule, population and horizon. Several
commercial tools expose quantities called policy value, reward, uplift or
expected utility. These quantities support different decisions. A static
treatment-assignment rule, a learning procedure that changes after each event,
a sequence of servicing actions, and an optimum within a supplied decision
model are distinct objects. The proposal should compare them at those
boundaries before deciding whether to buy a license, integrate an engine or
build a missing component.

The documented landscape already contains useful capabilities. Stata's
\texttt{cate} postestimation evaluates prespecified treatment-assignment
policies. The historical Azure Personalizer service documents offline
comparisons of learning policies. GeNIe and its SMILE engine provide decision
analysis and expected utilities in supplied influence diagrams. Amazon
Personalize documents held-out recommendation metrics. JMP's uplift platform
provides a separate targeting alternative. These are substantive comparators,
but they do not jointly establish a generally valid sequential causal-policy
service. The following comparison uses inspected official technical text;
none of these vendor engines was executed for this review.
\citep{exec13_ch_stata,exec16_cp_azure,exec16_cp_bayesfusion,exec16_cp_aws,exec13_ch_jmp_uplift}

\subsection{Four policy questions and one declared information clock}

For one binary treatment, let $X$ denote information available before treatment,
$Y(a)$ the potential outcome, and $\pi(x)\in[0,1]$ a prespecified probability
of assigning treatment. The static value is
\[
 V(\pi)=E\{\pi(X)Y(1)+[1-\pi(X)]Y(0)\}.
\]
This is the target displayed in Stata's policy-evaluation methods. The
command also compares two policies and accepts probability-valued policy
variables. Its synonym \texttt{estat tassigneval} should not be counted as
a separate method. The premises of the underlying treatment-effect analysis
still apply. A rule chosen by maximizing this same evaluation sample is a
different inferential target from a prespecified rule; a changing sequence of
actions is not represented merely by placing its last action in the treatment
column. \citep{exec13_ch_stata}

For sequential servicing, let $H_t$ be the observed history immediately before
action $A_t$, $R_t$ the subsequent net cash flow, $b_t(a\mid h)$ the logging
probability, and $\pi_t(a\mid h)$ the proposed policy. A finite-horizon target is
\[
 V_T(\pi)=E_\pi\left[\sum_{t=0}^{T-1}\gamma^tR_t\right],\qquad
 W_t=\prod_{s=0}^t\frac{\pi_s(A_s\mid H_s)}{b_s(A_s\mid H_s)}.
\]
The change-of-law argument needs consistent actions and outcomes, sufficient
histories for sequential assignment, positive logging probability wherever
the target policy acts, and an appropriate initial population and observation
law. Under a correctly specified observed-state experiment, the likelihood
ratios change the action distribution while leaving the shared environmental
mechanisms unchanged. Jiang and Li develop trajectory and stepwise importance
sampling and a sequential doubly robust construction; their discussion also
requires separating a data-derived target policy from the sample used for
evaluation. That result does not grant causal validity to arbitrary hidden
servicing states or misrecorded logging probabilities. \citep{jiang2016}

A supplied-model decision engine answers a third question: what action has
the highest expected utility if its probability tables, information sets and
utilities are correct? A recommendation metric answers a fourth: how well
does a learned ranking reproduce or predict held-out interactions? Either
can be valuable without identifying the effect of deploying a new servicing
policy. The commercial feature matrix should retain the target, source of
mechanisms, sampling unit and uncertainty claim for each quantity.

\subsection{A final-action adjustment can reverse a servicing decision}

Consider a generated two-stage law with $\gamma=1$. The first action $A_0$ is a fair binary
draw; the next observed state is $L_1=A_0$. The logger chooses $A_1=1$ with
probability $9/10$ after $A_0=0$ and $1/10$ after $A_0=1$. Define a terminal
net reward, in stipulated dollars, by
\[
 R=20A_0+100A_1-120A_0A_1-5(A_0+A_1).
\]
This law deliberately models an adverse interaction between repeated actions;
it is not an estimated banking mechanism. Its four paths are
\[
 \begin{array}{c|rrrr}
 (A_0,A_1)&(0,0)&(0,1)&(1,0)&(1,1)\\\hline
 P_b&1/20&9/20&9/20&1/20\\
 R&0&95&15&-10
 \end{array}
\]
For the policy setting both actions to one, direct substitution gives
$V=-10$. The full trajectory weight is twenty on path $(1,1)$ and zero
otherwise, hence $E_b[WR]=(1/20)(20)(-10)=-10$.
\citep{exec16_cp_exact}

Using only the history-conditional last-action ratio answers another question:
\[
 E_b\left[\frac{\mathbf1\{A_1=1\}}{b_1(1\mid A_0)}R\right]
 =\tfrac12(95)+\tfrac12(-10)=42.5.
\]
It retains the logger's first-action distribution. Using the marginal
last-action probability $P_b(A_1=1)=1/2$ instead gives
$2\{(9/20)95+(1/20)(-10)\}=84.5$. Neither positive number evaluates the
two-action target whose value is negative. The example specifies exactly
what a static-policy comparator can and cannot replace. More observations
of the same law would estimate the wrong target more precisely.
\citep{exec16_cp_exact}

There is also a precision cost even when the weighting is correct. Here
$E_bW=1$ and $E_bW^2=20$, so the population ratio
$(E_bW)^2/E_bW^2$ is $1/20$. Only one twentieth of trajectories follow
the target path in expectation. This ratio explains a loss of information;
it is not a finite-sample coverage guarantee or a universal effective sample
size. The number of logged events alone cannot settle whether a policy
comparison is precise enough for an investment decision.
\citep{exec16_cp_exact}

\subsection{Azure Personalizer is a historical workflow reference}

The acquired Microsoft pages describe Rank requests containing available
actions, action features, context and an event identifier, followed by a Reward
call. The inspected overview uses a reward between zero and one and permits
delayed feedback. It describes learning across users and expressly says that
the service does not maintain individual user histories. These are useful
workflow features, but a sequence of event-level calls is not itself a
specified customer-trajectory causal model. A monetary servicing objective
also needs an explicit reward transformation, treatment costs and horizon;
a click score cannot be silently substituted for discounted net recovery.
\citep{exec16_cp_azure}

The offline-evaluation pages compare the current online learning policy,
the application's default action, a random-action policy, uploaded alternatives
and policies found by optimization discovery. They display estimated average
rewards and confidence intervals, allow a policy to be downloaded or applied,
and recommend representative traffic and at least 50,000 events as a practical
starting point. The conceptual description replays chronological events with
a blank model for each learning policy and retains observed rewards when the
chosen action matches. A learning-policy configuration in this description
controls a changing learner; it is not necessarily a single fixed action rule.
\citep{exec16_cp_azure}

That conceptual account is insufficient to reconstruct a general estimator
or its uncertainty theorem. For an exact illustration, suppose two contexts
are equally common, a target always chooses action one, and its mean rewards
are 10 and 100. Let the logger's action-one probabilities be $0.9$ and $0.1$.
The target value is 55, but the mean among matching logged actions is
\[
 \frac{0.9(10)+0.1(100)}{0.9+0.1}=19.
\]
Weighting by the correct context-specific assignment probabilities recovers
55 in this population. This is a counterexample to matching alone under
nonuniform logging, not a claim that Microsoft's actual implementation uses
that incomplete estimator. The inspected pages do not supply all the
algorithmic and confidence-bound details needed to make such a runtime
claim. A successor comparison must acquire those details and verify the
logging distribution, retained-event denominator, replay updates and
policy-selection rule. \citep{exec16_cp_azure,exec16_cp_exact}

Availability changes the acquisition recommendation. The official pages state
that new resources became unavailable on 20 September 2023 and that the service
would retire on 25 August 2026. That retirement date precedes the proposal's
7 October 2026 cutoff. The pages recommend the open-source
\texttt{microsoft/learning-loop} project. Thus Personalizer is a historical
workflow comparator rather than a proposed new hosted dependency. The named
successor needs its own pinned code, rights and estimator review; the migration
link does not establish feature or statistical parity. The pages were acquired
after the cutoff and are labeled as archived documentation, with their
displayed update dates and source metadata retained. \citep{exec16_cp_azure}

\subsection{Recommendation quality, incremental value and capacity}

Amazon Personalize's inspected technical page defines train/test constructions,
metrics such as precision, recall and normalized discounted cumulative gain,
and a metric named \texttt{average\_rewards\_at\_k}. The latter normalizes
rewards from interactions with selected recommendations relative to rewards
from interactions with recommendations. Its name does not make it an estimate
of the population intervention contrast. The page also warns against comparing
solution metrics across different training data, and provides CLI and SDK
routes to retrieve metrics. Those are meaningful operational capabilities to
compare, while their documented target remains distinct from a causal policy
effect. This mutable page has no independently authenticated cutoff byte
identity; the source limitation is retained. \citep{exec16_cp_aws}

A generated capacity example makes the distinction operational. Suppose
80\% of customers have baseline recovery probability $0.9$ and a contact
effect $0.02$, while 20\% have baseline probability $0.1$ and effect $0.2$.
At USD~100 per recovery and USD~5 per contact, incremental net values are
$-3$ and $15$. Treating everyone gives only
$0.8(-3)+0.2(15)=0.6$ dollars per customer; contacting the high-increment
group gives $0.2(15)=3$. If capacity permits contacting 10\% of all customers,
randomly selecting half of that group yields $1.5$ dollars per customer.
Ranking by treated response probability instead favors the low-increment
group, whose response probability is $0.92$ rather than $0.3$, and gives
$0.1(-3)=-0.3$ dollars per customer at the same capacity.
\citep{exec16_cp_exact}

The values are stipulated, but the comparison identifies a product requirement:
retain the baseline outcome, incremental effect, cost, assignment probability
and capacity separately. JMP's inspected uplift platform explicitly chooses
partition splits to emphasize treatment differences, using treatment-by-split
interactions for continuous or categorical responses. It is consequently a
relevant targeting comparator. Its older guide's validation and treatment-level
details are scoped to that guide; the current overview cannot establish that
every edition implements the same procedure. Stata's prespecified policy-value
target provides another relevant comparison once the net outcome and population
are aligned. Neither source licenses an automatic inference from a fitted
ranking to an independently validated policy gain. \citep{exec13_ch_jmp_uplift,exec13_ch_stata}

\subsection{Selecting a policy is an additional experiment}

Suppose $m$ candidate policies all have true incremental value zero, and each
evaluation has an independent error equal to $-4$ or $4$ with equal probability.
Every individual estimate is unbiased. Reporting the largest estimate gives
\[
 E\left[\max_{j\leq m}\widehat V_j\right]
 =4\{1-2^{1-m}\}.
\]
The maximum is $-4$ only when all errors are negative and is $4$ otherwise.
Its expectation is zero for one candidate, two for two candidates, and
$127/32$, approximately $3.969$, for eight. A fresh independent mean-zero
evaluation of the selected index has expectation zero. This exact finite
example explains the role of an independent final evaluation; it is not a
proposed universal subtraction rule for a commercial optimizer.
\citep{exec16_cp_exact}

The candidate list, training and tuning records, validation sample and final
policy identity should therefore be part of the comparison. If a product
supports repeated online adaptation, the evaluation must reproduce that
adaptation under an appropriate experiment rather than call it a frozen-policy
test. The best policy among a declared finite list also differs from an
unknown optimum over all feasible strategies. The proposed reporting should
state which value is being estimated and preserve the selection uncertainty
or the independence that justifies its exclusion. \citep{jiang2016,exec16_cp_matrix}

\subsection{Influence diagrams support a different reuse decision}

SMILE's inspected manual represents controllable decisions as nodes with
available states, value nodes with utility tables, and multi-attribute utility
nodes with weights or expressions. It exposes indexing parents, including
earlier decisions and chance observations relevant before a decision. After
inference, its matrices contain policy-indexed probabilities or expected
utilities, depending on node type. GeNIe supplies the associated decision
framing, preference elicitation and value-of-information workflow. These
capabilities are relevant to a bank that can supply a credible finite decision
model. \citep{exec16_cp_bayesfusion}

The proposed reuse choice is a licensed supplied-model comparator. A bridge
must preserve which observations precede which actions, decision-state order,
utility units, the interpretation of indexing parents, and the distinction
between evaluating a supplied policy and optimizing within the model.
Maximizing a supplied utility table cannot validate its transition probabilities
against a selectively logged bank history. Nor does the presence of hidden
chance nodes identify a partially observed causal system from arbitrary data.
The inherited SMILE licensing documentation remains relevant; public manuals
do not grant embedding rights or establish a current price. Porting the vendor
engine is not the proposed route. \citep{exec16_cp_bayesfusion,exec13_smile_license}

\subsection{What the proposed product still needs}

The common local reference revision examined in this probe exposed six
bounded scientific operations; its declared operation schema and estimator
dispatch did not provide a general policy-evaluation service. Three actual local requests naming fixed-history,
learned-policy validation and static-assignment value methods returned
\texttt{unsupported} with no result. This observation is restricted to that
service. The repository also contains separate history-evaluation and
controlled-model research modules; their existence and selected demonstrations
are not equivalent to a qualified common library/CLI/MCP policy interface.
The later fitted servicing and selected-policy studies
(Chapters~\ref{gap29:servicing} and~\ref{gap29:policy-uncertainty}) develop
bounded inferential evidence. The fitted-summary interface demonstration in
Chapter~\ref{gap29:public} is retained separately. None establishes a general
raw-history policy estimator or validates bank logging probabilities.
\citep{exec16_cp_runtime}

The concrete proposal gap is therefore larger than adding another estimator
identifier. A policy request must bind its initial population, physical and
availability clocks, full history variables, action probabilities, observation
and censoring process, discounting, net-reward map, selection/training identity
and final evaluation units. Its response should distinguish a static rule,
a fixed sequence policy, an adapting learner and a supplied-model optimum.
It must retain support failures and uncertainty limits. The documented
commercial workflow features---policy comparison, model export, diagnostics
and decision-model inspection---are useful requirements for this design, while
the scientific meaning is checked independently. \citep{exec16_cp_matrix}

The bounded portfolio covers eleven named families. Stata and JMP support
documented static alternatives; Azure supplies a retired workflow reference;
Amazon supplies a recommendation-evaluation comparator; and BayesFusion
supplies model-based decision analysis. MATLAB's selected official routes were
inaccessible. SAS reinforcement-learning routes and HUGIN's selected routes
did not yield technical originals. The inherited SAS survival, IBM regression,
Mplus dynamic-model and Bayesia causal-model sources cannot settle their
policy-specific capabilities. These entries remain unresolved with concrete
manual/version acquisition routes. No absence of a policy capability is
inferred from an unavailable page. Exact editions, entitlements, prices and
runtime behavior remain open across the prospective commercial trials.
\citep{exec16_cp_matrix}

An economical prototype can begin with the exact mechanisms above and the
existing open-source reference comparisons. Public data can supply realistic
covariates, event clocks and missingness patterns, while generated assignments
and outcomes provide declared causal truth. Compare one static policy,
one history policy and one supplied-model decision on matched eligible
questions; retain a predictive recommendation baseline as an explicitly
different alternative. Freeze final trajectories, candidate policies, metrics,
cost accounting and stopping rules before evaluation. Native commercial
execution follows only for a selected edition with available rights. Bank
experiments become relevant when the prototype identifies a consequential
logging, support, hidden-state or transport uncertainty that public/generated
evidence cannot settle. This route makes the proposal executable while leaving
the unresolved general research questions visible. \citep{exec16_cp_matrix}
\label{exec16:commercial-policy-end}
